\documentclass[12pt,a4paper]{article}
\usepackage[UTF8]{ctex}
\usepackage{amsmath,amssymb}
\usepackage{graphicx}
\usepackage{booktabs}
\usepackage{geometry}
\usepackage{siunitx}
\usepackage{ulem}
\geometry{left=2.5cm,right=2.5cm,top=2.5cm,bottom=2.5cm}

\title{\Huge{\textbf{\uuline{核\ 磁\ 共\ 振\ 实\ 验\ 报\ 告}}}}
\author{\uline{24级少院} \quad 姓名: \uline{徐子钦} \quad 学号: \uline{PB24000224}}
\date{2026.5.1}

\begin{document}
\maketitle

\begin{abstract}
本实验以核磁共振(NMR)原理为基础，通过在恒定磁场中施加射频脉冲，观测样品的自由感应衰减(FID)与自旋回波信号，测定核磁共振频率及弛豫时间 $T_1$ 与 $T_2$。本报告依据实验手册（核磁共振.pdf）整理实验原理与方法，数据部分留空，待实验测量后填入并完成数据处理与讨论。
\end{abstract}

\section{实验题目}
核磁共振（NMR）实验

\section{实验目的}
\begin{enumerate}
    \item 理解核磁共振的基本物理机制与共振条件。
    \item 掌握使用脉冲射频激发与检测自由感应信号（FID）和回波信号的方法。
    \item 测定样品的共振频率及弛豫时间 $T_1$、$T_2$，并进行误差分析。
\end{enumerate}

\section{实验原理}
\subsection{自旋与磁矩}
原子核具有自旋角动量 $\mathbf{I}$ 与磁矩 $\boldsymbol{\mu}=\gamma \mathbf{I}$，其中 $\gamma$ 为回旋比。置于外加静磁场 $\mathbf{B}_0$ 中，核自旋能级发生塞曼分裂。

\subsection{拉莫尔进动与共振条件}
核磁矩在 $\mathbf{B}_0$ 中发生拉莫尔进动，角频率为
$$\omega_0 = \gamma B_0,$$
对应频率 $f_0 = \omega_0/2\pi$。当施加频率为 $f_0$ 的横向射频场 $B_1$ 时发生能量吸收，产生可检测的电磁信号。

\subsection{弛豫过程}
横纵弛豫分别用时间常数 $T_2$（横向）与 $T_1$（纵向）描述。弛豫动力学可用布洛赫方程描述，自由感应衰减（FID）与回波信号的幅度随时间按指数规律衰减，典型形式为
$$M_x(t)\propto e^{-t/T_2},\qquad M_z(t)\propto 1-2e^{-t/T_1}\ \ (\text{视具体脉冲序列而定}).$$

\subsection{测量方法概述}
常用测量 $T_1$ 的方法为反转恢复（inversion-recovery）；测量 $T_2$ 的常用方法为自旋回波（spin-echo）或Carr–Purcell–Meiboom–Gill (CPMG) 序列。实验室通过调整脉冲角度（90°/180°）、脉间延时 $\tau$ 并拟合回波幅度或纵向恢复曲线获得弛豫时间。

\section{实验器材}
\begin{itemize}
    \item 样品与样品架（注明样品种类）；
    \item 恒定磁场主磁体（标明 $B_0$ 值或测量方法）；
    \item 脉冲发射与接收系统（功率放大器、射频线圈）；
    \item 示波器或数字接收器用于记录 FID/回波波形；
    \item 计算机及数据处理软件。
\end{itemize}

\section{实验步骤（概括）}
\begin{enumerate}
    \item 将样品安装到射频线圈中心，固定位置。
    \item 打开磁体、调节匀场（必要时做shimming），记录主磁场 $B_0$。
    \item 设置射频频率在预计共振附近，扫描频率以确定共振峰 $f_0$。
    \item 确定 90° 与 180° 脉冲宽度（通过脉冲归一化或方位调谐）。
    \item 采用 FID、反转恢复与自旋回波序列分别测量共振信号与弛豫过程，记录原始波形数据。
    \item 将测得的时域信号进行指数拟合或逆变换以求取 $T_1$, $T_2$ 及谱线宽度等参数。
\end{enumerate}

\section{原始数据（待填写）}
\subsection*{实验参数}
\begin{center}
\begin{tabular}{ll}
\toprule
项目 & 值 \\
\midrule
主磁场 $B_0$ &  \\
共振频率 $f_0$ &  \\
样品名称/编号 &  \\
采样率/带宽 &  \\
脉冲宽度（90°/180°） &  \\
测量温度 &  \\
\bottomrule
\end{tabular}
\end{center}

\subsection*{共振频率扫描数据}
（在此处粘贴或附上扫描数据表格或图像，占位：\includegraphics[width=0.8\textwidth]{resonance_scan.png}）

\subsection*{FID 与回波测量原始数据}
\begin{itemize}
    \item FID 原始波形：占位 \includegraphics[width=0.8\textwidth]{FID_placeholder.png}
    \item 回波序列原始波形：占位 \includegraphics[width=0.8\textwidth]{echo_placeholder.png}
\end{itemize}

\subsection*{弛豫时间测量表（空）}
\begin{center}
\begin{tabular}{cccc}
\toprule
序号 & 延时 $\tau$ (s) & 回波幅度 (V) & 备注 \\
\midrule
 & & & \\
 & & & \\
 & & & \\
\bottomrule
\end{tabular}
\end{center}

\section{数据处理与计算（方法与公式）}
\subsection{共振频率与回旋比}
由测得共振频率 $f_0$ 与主场 $B_0$ 可求回旋比
$$\gamma = \frac{2\pi f_0}{B_0}.$$ 
若样品为已知核种（如质子），可用此公式检验 $B_0$ 的标定。

\subsection{弛豫时间拟合}
\paragraph{T_2:} 自旋回波幅度 $A(\tau)$ 随延时衰减，一般按指数拟合
$$A(\tau)=A_0\,e^{-2\tau/T_2},$$
或根据所用序列的具体时间因子调整拟合公式。
\paragraph{T_1:} 反转恢复序列的纵向磁化强度 $M_z(t)$ 可按
$$M_z(t)=M_0\left(1-2e^{-t/T_1}\right)$$
或相应的补正规则拟合以求 $T_1$。

\subsection{不确定度估计（示例）}
对拟合参数的不确定度可由最小二乘拟合的协方差矩阵给出；以 $T_2$ 为例，若拟合得 $\hat{T}_2$ 与标准误差 $u(T_2)$，则给出不确定度表达为 $T_2=\hat{T}_2\pm u(T_2)$（适当时给出置信区间，如 $k=2$）。

\section{结果与讨论（占位）}
在完成测量并填写原始数据后，按上节所述方法计算 $f_0$, $\gamma$, $T_1$, $T_2$ 并在此处给出数值结果、不确定度及讨论。讨论要点建议包括：谱线宽度与 $T_2$ 的关系、样品成分或杂质对 $T_1,T_2$ 的影响、系统噪声与测量精度、磁场不均匀性对结果的影响、以及与文献值或手册预期值的比较。

\section{结论（占位）}
在填写并处理完数据后，在此处给出实验的主要结论，例如测得的 $T_1,T_2$ 值、共振频率与理论/标定值的一致性评价，以及实验局限与改进建议。

\section{参考文献}
（请在此处列出实验手册与参考文献）

\section*{附录：原始数据文件清单}
\begin{itemize}
    \item resonance_scan.csv（共振扫描原始数据）
    \item FID_waveforms/（原始时域波形文件）
    \item echo_waveforms/（回波序列原始文件）
\end{itemize}

\end{document}
